
\subsubsection{4.1 Dataset}

Waterbirds \citep{sagawa2020groupdro} contains 11,788 images from the Caltech-UCSD Birds-200-2011 dataset placed on backgrounds from the Places dataset. The dataset has 4,795 training images, 1,199 validation images, and 5,794 test images. The 4 spurious groups have the following training distribution:

\begin{itemize}
\item Group 0 (landbird, land): 3,498 samples (73\%)
\item Group 1 (landbird, water): 184 samples (4\%) [minority]
\item Group 2 (waterbird, water): 980 samples (20\%)
\item Group 3 (waterbird, land): 133 samples (3\%) [minority]
\end{itemize}

The spurious correlation is 93\% in the training split.

\subsubsection{4.2 Implementation}

All experiments used PyTorch 1.13 and Python 3.9. Training was conducted on a single NVIDIA GPU. The SimCLR and LA-SSL implementations were validated through comprehensive unit and integration tests. The h-e1 implementation passed 43/43 tests with 100\% software development document (SDD) compliance (15/15 tasks). The h-m-integrated implementation passed 5/5 tests with 100\% SDD compliance (43/43 tasks).

Training hyperparameters were determined through grid search for learning rate $\in$ {0.01, 0.001, 0.0001} and weight decay $\in$ {$10^-4, 10^-5, 10^-6$}. Linear probe training used 20 epochs with batch size 32 and early stopping based on validation WGA.

\subsubsection{4.3 Proof-of-Concept Protocol}

Experiments were conducted as proof-of-concept validation with 20-epoch training runs. This duration was selected to validate implementation correctness and assess whether the mechanism hypotheses showed evidence of support before committing to full-scale 100-epoch experiments requiring 48-96 GPU hours. Checkpoints were saved at epochs 10 and 20 for analysis.
